%%%%%%%%%%%%%%%%%%%%%%%%%%%%%%%%%%%%%%%%%%%%%%%%%%%%%%%%%%%%%%%%%%%%%%
% Lớp 11 | Dãy số. Cấp số cộng. Cấp số nhân | Cấp số nhân
% Dạng 1D2?3-7: Kết hợp cấp số nhân và cấp số cộng
% Số câu: 23  (sinh tự động bởi tach_cau.py; từ nay thầy sửa trực tiếp file này)
%%%%%%%%%%%%%%%%%%%%%%%%%%%%%%%%%%%%%%%%%%%%%%%%%%%%%%%%%%%%%%%%%%%%%%

%%%%%%%%%%%%%%%% TRẮC NGHIỆM 4 PHƯƠNG ÁN %%%%%%%%%%%%%%%%

% nguồn: Ok_Toan11_TracNghiem_Lan2
\begin{ex}%[1D2H3-7]%[ma:8710dfaa]
	Trong các phát biểu sau, phát biểu nào là  \textbf{sai}?
	\choice 
		{Dãy số có tất cả các số hạng bằng nhau là một  cấp số nhân}
		{Dãy số có tất cả các số hạng bằng nhau  là một cấp số cộng}
		{Một cấp số cộng có công sai dương là  một dãy số tăng}
		{\True Một cấp số  cộng có công sai dương là một dãy số dương}
	 \loigiai{
		\begin{itemize}
			 \item
			Mệnh đề $\mathrm{A}$ đúng vì dãy  số đã cho là cấp số nhân với công bội $q=1$.
			\item 
			Mệnh đề $\mathrm{B}$ đúng vì  dãy số đã cho là cấp số cộng với công sai $d=0$.
			 \item
			Mệnh  đề $\mathrm{C}$ đúng vì dãy số đã cho là cấp số cộng có công sai dương nên $u_{n+1}-u_n=d>0 \Rightarrow u_{n+1}>u_n$.
			\item 
			Mệnh đề $\mathrm{D}$ sai. Ví dụ dãy $-5$; $-2$;  $1$; $4$; $\ldots$ là dãy số có số $d=3>0$ nhưng không phải là dãy số dương.
		\end{itemize}
	} 
\end{ex}

% nguồn: Ok_Toan11_TracNghiem_Lan2
\begin{ex}%[1D2H3-7]%[ma:0c7e088c]
	Cho dãy số $(u_n)$ có số hạng tổng quát là $u_n = 3\cdot 2^{n+1}$  $(\forall n \in \mathbb{N}^*)$. Chọn kết luận đúng
	\choice 
		{\True Dãy số là cấp  số nhân có số hạng đầu $u_1 = 12$}
		{Dãy số là cấp số cộng có công sai $d = 2$} 
		 {Dãy số là cấp số cộng có số hạng đầu $u_1 = 6$}
		{Dãy số là cấp số nhân có công  bội $q = 3$}
	 \loigiai{
		Dãy số $(u_n)$ có số hạng tổng quát là  $u_n = 3\cdot 2^{n+1} \Rightarrow u_{n+1} = 3\cdot 2^{n+2}$.\\
		 Xét thương $\dfrac{u_{n+1}}{u_n} = \dfrac{3\cdot 2^{n+2}}{3\cdot 2^{n+1}} = 2 = \text{const}$ với $\forall n \in \mathbb{N}^*$ nên $(u_n)$ là một cấp số nhân có công bội $q = 2$ và số hạng đầu là $u_1 = 3\cdot 2^{1+1} = 12$.
	 }
\end{ex}

% nguồn: Ok_Toan11_TracNghiem_Lan2
\begin{ex}%[1D2V3-7]%[ma:6db23003]
	Cho bốn số $a$, $b$, $c$, $d$ theo thứ tự đó tạo thành cấp số nhân với công  bội khác $1$. Biết tổng ba số hạng đầu bằng $\dfrac{148}{9}$, đồng thời theo thứ tự đó chúng lần lượt là số hạng thứ nhất, thứ tư và thứ tám của một cấp số cộng. Tính giá trị biểu thức $T=a-b+c-d$.
	\choice 
		{$T=\dfrac{101}{27}$} 
		 {$T=\dfrac{100}{27}$}
		{\True $T=-\dfrac{100}{27}$} 
		 {$T=-\dfrac{101}{27}$}
	\loigiai{ 
		Ta  có $\heva{&ac=b^2\quad(1)\\&bd=c^2 \quad(2)\\&a+b+c=\dfrac{148}{9}\quad(3).}$\\
		Và cấp số  cộng có $u_1=a$, $u_4=b$, $u_8=c$.\\
		Gọi $x$ là công  sai của cấp số cộng. Vì cấp số nhân có công bội khác $1$ nên $x\ne 0$.\\
		 Ta có $\heva{&b=a+3x\\&c=a+7x}\quad(4)$.\\
		Từ $(1)$ và  $(4)$ ta được $a(a+7x)=(a+3x)^2\Leftrightarrow ax-9x^2=0$.\\
		Do  $x\ne 0$ nên $a=9x$.\\
		Từ $(3)$ và $(4)$, suy ra $3a+10x=\dfrac{148}{9}$.\\ 
		 Do đó $\heva{&a=4\\&x=\dfrac{4}{9}}\Rightarrow \heva{&b=\dfrac{16}{3}\\&c=\dfrac{64}{9}\\&d=\dfrac{256}{27}.}$\\
		 Vậy $T=a-b+c-d=-\dfrac{100}{27}$.
	 }
\end{ex}

% nguồn: Ok_Toan11_TracNghiem_Lan2
\begin{ex}%[1D2V3-7]%[ma:bc7cb445]
	Ba số $x$, $y$, $z$ lập thành một cấp số cộng và có tổng bằng $21$. Nếu lần lượt thêm các số $2$; $3$; $9$ vào ba số đó thì  được ba số lập thành một cấp số nhân. Tính $F=x^2+y^2+z^2$.
	\choice 
		{$F=389$  hoặc $F=395$}
		{$F=395$ hoặc $F=179$} 
		{\True $F=389$ hoặc  $F=179$}
		{$F=441$  hoặc $F=357$}
	\loigiai{ 
		Theo tính chất của cấp số cộng,  ta có $x+z=2y$.\\
		Kết  hợp với giả thiết $x+y+z=21$, ta suy ra $3y=21\Leftrightarrow y=7$.\\
		Gọi $d$ là công sai của cấp số cộng thì $x=y-d=7-d$ và $z=y+d=7+d$.\\ 
		Sau khi thêm các số $2$; $3$; $9$ vào ba số $x$, $y$, $z$ ta được ba  số là $x+2$, $y+3$, $z+9$ hay $9-d$, $10$, $16+d$.\\
		Theo tính chất  của cấp số nhân, ta có $\left(9-d\right)\left(16+d\right)=10^2\Leftrightarrow d^2+7d-44=0$.\\
		Giải phương trình  ta được $d=-11$ hoặc $d=4$.\\
		\begin{itemize}
			\item Với $d=-11$,  cấp số cộng $18$, $7$, $-4$. Lúc này $F=389$.
			\item Với $d=4$, cấp số cộng $3$, $7$, $11$. Lúc  này $F=179$.
		\end{itemize}
	} 
\end{ex}

% nguồn: Ok_Toan11_TracNghiem_Lan2
\begin{ex}%[1D2V3-7]%[ma:d2eda509]
	Cho dãy số $\left(a_n \right)$ xác định bởi $a_1=5, a_{n+1}=q.a_n+3$ với mọi $n\ge 1$, trong đó $q$ là hằng số, $q\ne 0$, $q\ne 1$. Biết công thức số hạng tổng  quát của dãy số viết được dưới dạng $a_n=\alpha.q^{n-1}+\beta \dfrac{1-q^{n-1}}{1-q}$. Tính $\alpha+2\beta$?
	 \choice
		 {$13$}
		{$9$} 
		{\True  $11$}
		{$16$} 
	 \loigiai{
		\begin{itemize}
			 \item Cách 1.\\
			Ta  có $a_{n+1}-k=q(a_n-k)\Leftrightarrow k-kq=3\Leftrightarrow k=\dfrac{3}{1-q}$.\\
			Đặt $v_n=a_n-k\Rightarrow  v_{n+1}=q\cdot v_n=q^2\cdot v_{n-1}=\ldots=q^n\cdot v_1$.\\
			Khi đó $v_n=q^{n-1}\cdot v_1=q^{n-1}\cdot  \left(a_1-k\right)=q^{n-1}\cdot \left(5-\dfrac{3}{1-q} \right)$.\\
			Vậy $a_n=v_n+k=q^{n-1}\cdot \left(5-\dfrac{3}{1-q} \right)+k=q^{n-1}\cdot  \left(5-\dfrac{3}{1-q} \right)+\dfrac{3}{1-q}=5\cdot q^{n-1}+3\cdot \dfrac{1-q^{n-1}}{1-q}$.\\
			Do  đó $\alpha=5$; $\beta=3 \Rightarrow \alpha+2\beta=5+2\cdot 3=11$.
			\item Cách  2.\\
			Theo giả  thiết ta có $a_1=5$, $a_2=5q+3$.\\
			 Áp dụng công thức tổng quát, ta được $\heva{&a_1=\alpha\cdot q^{1-1}+\beta \dfrac{1-q^{1-1}}{1-q}=\alpha\\&a_2=\alpha\cdot q^{2-1}+\beta \dfrac{1-q^{2-1}}{1-q}=\alpha q+\beta.}$\\
			 Suy ra $\heva{&5=\alpha\\&5q+3=\alpha q+\beta}$, hay $\heva{&\alpha=5\\&\beta=3.}$
			Do  đó $\alpha+2\beta=5+2\cdot 3=11$.
		\end{itemize}
	 }
\end{ex}

% nguồn: Ok_Toan11_TracNghiem_Lan2
\begin{ex}%[1D2V3-7]%[ma:4e4f67db]
	Các số $x+6y$, $5x+2y$, $8x+y$ theo thứ tự đó lập thành một cấp số cộng, đồng thời, các số $x+\dfrac{5}{3}$, $y-1$, $2x-3y$  theo thứ tự đó lập thành một cấp số nhân. Hãy tìm $x$ và $y$.
	\choice 
		{\True $x=-3$, $y=-1$ hoặc $x=\dfrac{3}{8}$, $ y=\dfrac{1}{8}$} 
		{$x=3$, $y=1$ hoặc $x=-\dfrac{3}{8}$, $  y=-\dfrac{1}{8}$}
		 {$x=24$, $y=8$ hoặc $x=-3$, $y=-1$}
		{$x=-24$, $y=-8$ hoặc $x=3$,  $y=1$}
	\loigiai{ 
		\begin{itemize}
			\item Ba số $x+6y$, $5x+2y$, $8x+y$  lập thành cấp số cộng nên $\left(x+6y\right)+\left(8x+y\right)=2\left(5x+2y\right)\Leftrightarrow x=3y$.
			\item Ba số $x+\dfrac{5}{3}$, $ y-1$, $2x-3y$ lập thành cấp số nhân  nên $\left(x+\dfrac{5}{3} \right)\left(2x-3y\right)=\left(y-1\right)^2$.
		\end{itemize}
		Thay $x=3y$ vào ta được $8y^2+7y-1=0\Leftrightarrow y=-1$ hoặc $y=\dfrac{1}{8}$.\\ 
		Với $y=-1$ thì  $x=-3$; với $y=\dfrac{1}{8}$ thì $x=\dfrac{3}{8}$.
	 }
\end{ex}

% nguồn: Ok_Toan11_TracNghiem_Lan2
\begin{ex}%[1D2V3-7]%[ma:c395090f]
	Cho ba số thực $x$, $y$, $z$ trong đó $x\ne 0$. Biết rằng $x$, $2y$, $3z$ lập  thành cấp số cộng và $x$, $y$, $z$ lập thành cấp số nhân. Tìm công bội $q$ của cấp số nhân đó.
	 \choice
		{\True $\hoac{  &q=1 \\ &q=\dfrac{1}{3} }$}
		{$\hoac{  &q=\dfrac{1}{3} \\ &q=\dfrac{2}{3} }$}
		{$q=2$} 
		{$q=1$} 
	 \loigiai{
	$x,y,z$ lập thành cấp số nhân có công bội $q$ nên  $y=qx$, $z=q^2x$.\\
	$x,2y,3z$ lập thành  cấp số cộng nên $2y=\dfrac{2+3z}{2}\Rightarrow 2qx=\dfrac{x+3q^2x}{2}$.\\
	Vì $x\ne 0$  nên $2qx=\dfrac{x+3q^2x}{2}\Rightarrow 4q=1+3q^2\Rightarrow \hoac{ &q=1 \\ &q=\dfrac{1}{3}.}$
	 }
\end{ex}

% nguồn: Ok_Toan11_TracNghiem_Lan2
\begin{ex}%[1D2V3-7]%[ma:35d4bbbc]
	Các số $x+6y$,  $5x+2y$, $8x+y$ theo thứ tự đó lập thành một cấp số cộng; đồng thời các số $x-1$, $y+2$, $ x-3y$ theo thứ tự đó lập thành một cấp số nhân. Tính $x^2+y^2$.
	 \choice
		{\True  $x^2+y^2=40$}
		{$x^2+y^2=25$} 
		{$x^2+y^2=100$} 
		{$x^2+y^2=10$} 
	 \loigiai{
		 Theo giả thiết ta có
		\begin{align*}
			 &\heva{&{\left(x+6y\right)+\left(8x+y\right)=2\left(5x+2y\right)} \\
			 &{\left(x-1\right)\left(x-3y\right)=\left(y+2\right)^2}
			 }\\
			\Leftrightarrow&  \heva{&{x=3y} \\
			&{\left(3y-1\right)\left(3y-3y\right)=\left(y+2\right)^2} 
			 } \Leftrightarrow \heva{&{x=3y} \\
			&{0=\left(y+2\right)^2} 
			 } \Leftrightarrow \heva{&{x=-6} \\
			 &{y=-2}.
			 }
		\end{align*}
		Suy ra $x^2+y^2=40$. 
	} 
\end{ex}

% nguồn: Ok_Toan11_TracNghiem_Lan2
\begin{ex}%[1D2V3-7]%[ma:e9150947]
	Ba số $x$; $y$; $z$ theo  thứ tự lập thành một cấp số nhân với công bội $q$ khác $1$; đồng thời các số $x$; $2y$; $3z$ theo thứ tự lập thành một cấp số cộng với công sai khác $0$. Tìm giá trị của $q$.
	\choice 
		{\True  $q=\dfrac{1}{3}$}
		{$q=\dfrac{1}{9}$} 
		{$q=-\dfrac{1}{3}$} 
		{$q=-3$} 
	 \loigiai{
		Ta  có
		\[\heva{&{y=xq; z=xq^2}  \\
		&{x+3z=2(2y)} 
		} \Rightarrow x+3xq^2=4xq\Rightarrow  x\left(3q^2-4q+1\right)=0\Leftrightarrow \hoac{&{x=0} \\
		&{3q^2-4q+1=0}. 
		}\] 
		Nếu $x=0\Rightarrow y=z=0\Rightarrow$ công sai của cấp số cộng: $x$; $2y$; $3z$ bằng  $0$.\\
		Nếu $3q^2-4q+1=0\Leftrightarrow \hoac{&{q=1} \\ 
		 &{q=\dfrac{1}{3}}
		 } \Leftrightarrow q=\dfrac{1}{3} \;\left(q\not=1\right)$.
	} 
\end{ex}

% nguồn: Ok_Toan11_TracNghiem_Lan2
\begin{ex}%[1D2V3-7]%[ma:0bf2bb14]
	Ba số theo thứ tự lập thành một cấp số nhân có số hạng cuối lớn hơn số hạng đầu $16$ đơn vị. Ba số đó là các số hạng thứ nhất, thứ hai và thứ năm của một cấp số cộng. Tìm ba số đó. 
	\choice 
		 {$2$, $6$, $18$}
		{$4$, $8$,  $20$}
		{$\dfrac{1}{3}$,  $\dfrac{7}{3}$, $\dfrac{49}{3}$}
		{$4$,  $4\sqrt{5}$, $20$}
	 \loigiai{
		Ta gọi ba số  đó lần lượt là $a,b,c$ và $d$ là công sai của cấp số cộng.\\
		Theo đề  bài ta có $\heva{&c=a+16\\&c=a+4d}\Rightarrow d=4$.
		Ngoài ra $b^2=ac\Leftrightarrow  (a+4)^2=a(a+16)\Leftrightarrow a=2$.\\
		Suy ra $b=6$,  $c=18$.\\
		Vậy các  số cần tìm là $2$, $6$, $18$.
	} 
\end{ex}

% nguồn: Ok_Toan11_TracNghiem_Lan2
\begin{ex}%[1D2V3-7]%[ma:6cdbbe8e]
	Ba số $x$, $y$, $z$ theo thứ tự lập thành một cấp số cộng tăng có tổng bằng $24$. Nếu cộng thêm lần lượt các  số $1$, $4$, $13$ vào ba số $x$, $y$, $z$ ta được ba số theo thứ tự lập thành cấp số nhân. Tính giá trị biểu thức $P=x^2+y^2+z^2$.
	 \choice
		{$200$} 
		 {\True $210$}
		 {$220$}
		 {$190$}
	\loigiai{ 
		Ba số $x$, $y$, $z$ theo thứ tự lập thành  một cấp số cộng có tổng bằng $24$ nên ta có hệ phương trình
		 \[\heva{&x+y+z=24\\&x+z=2y} \Rightarrow 3y=24\Rightarrow y=8.\]
		Ta viết lại $3$ số $x$, $y$, $z$  lần lượt bằng $8-d$, $8$, $8+d$.\\
		Nếu cộng thêm lần lượt các số  $1$, $4$, $13$ vào ba số $x$, $y$, $z$ ta được ba số là $9-d$, $12$, $21+d$.\\
		Vì ba số này theo  thứ tự lập thành cấp số nhân nên ta có phương trình
		\begin{eqnarray*}
			 (9-d)(21+d)=12^2 &\Leftrightarrow & d^2+12d-45=0\\
			&\Leftrightarrow  & \hoac{&d=3\\&d=-15.}
		\end{eqnarray*}
		Vì cấp số cộng tăng nên $d> 0\Rightarrow d=3\Rightarrow$ ba số $x$, $y$, $z$ lần lượt bằng $5$, $8$,  $11$.\\
		 Suy ra $P=x^2+y^2+z^2=5^2+8^2+11^2=210$.
	} 
\end{ex}

% nguồn: Ok_Toan11_TracNghiem_Lan2
\begin{ex}%[1D2V3-7]%[ma:1a310c36]
	Cho bố số $a$, $b$, $c$, $d$ biết rằng $a$, $b$, $c$ theo thứ tự đó lập thành một cấp số nhân công bội $q > 1$; còn $b$, $c$, $d$ theo thứ tự đó lập thành cấp số cộng. Tìm $q$ biết rằng $a+d=14$ và $b+c=12$. 
	\choice 
		{$q=\dfrac{18+\sqrt{73}}{24}$} 
		 {\True $q=\dfrac{19+\sqrt{73}}{24}$}
		{$q=\dfrac{20+\sqrt{73}}{24}$} 
		 {$q=\dfrac{21+\sqrt{73}}{24}$}
	\loigiai{ 
		Giả sử $a$, $b$, $c$ lập thành  cấp số cộng công bội $q$. Khi đó theo giả thiết ta có
		\[\heva{& {b=aq, c=aq^2} \\ &{b+d=2c} \\ &{a+d=14} \\ &{b+c=12} } \Rightarrow \heva{ &{aq+d=2aq^2\;\;\qquad(1)} \\ &{a+d=14\qquad\qquad(2)} \\ &{a\left(q+q^2\right)=12\qquad(3)}  }\]
		Khi $q>1$, từ $(2)$ và $(3)$ ta có $d=14-a$, $a=\dfrac{12}{q+q^2}$ thay vào $(1)$  ta được
		\begin{align*}
			\dfrac{12q}{q+q^2}+\dfrac{14q^2+14q-12}{q+q^2}=\dfrac{24q^3}{q+q^2} &\Leftrightarrow 12q^3-7q^2-13q+6=0  \\ &\Leftrightarrow \left(q+1\right)\left(12q^2-19q+6\right)=0\\
			 &\Leftrightarrow q=\dfrac{19\pm \sqrt{73}}{24}.
		\end{align*}
		Vì $q >  1$ nên $q=\dfrac{19+\sqrt{73}}{24}$.
	} 
\end{ex}

% nguồn: Ok_Toan11_TracNghiem_Lan2
\begin{ex}%[1D2V3-7]%[ma:6077725a]
	Cho ba số $a$, $b$, $c$  là ba số liên tiếp của một cấp số cộng có công sai là $2$. Nếu tăng số thứ nhất thêm $1$, tăng số thứ hai thêm $1$ và tăng số thứ ba thêm $3$ thì được ba số mới là ba số liên tiếp của một cấp số nhân. Tính $a+b+c$.
	 \choice
		 {$12$}
		 {$18$}
		 {$3$}
		{\True $9$} 
	 \loigiai{
		$a$, $b$, $c$ là ba số hạng  liên tiếp của cấp số cộng có công sai bằng $d=2\Rightarrow \heva{&b=a+2\\&c=a+4.}$\\
		Ba số $a+1$, $a+3$, $a+7$ là ba số hạng liên tiếp  của một cấp số nhân
		\[\Leftrightarrow (a+3)^2=(a+1)(a+7) \Leftrightarrow a^2+6a+9=a^2+8a+7\Leftrightarrow  2a=2\Leftrightarrow a=1.\]
		Do  đó $a+b+c=3a+6=9$.
	} 
\end{ex}

% nguồn: Ok_Toan11_TracNghiem_Lan2
\begin{ex}%[1D2V3-7]%[ma:f34e9c81]
	Ba số dương $x$, $y$, $z$ theo thứ tự lập thành một cấp số cộng và có tổng bằng $30$. Biết  $x+2$; $y+2$; $z+18$ theo thứ tự lập thành một cấp số nhân. Tính $T=x^2+z^2$.
	 \choice
		{\True  $T=328$}
		{$T=424$} 
		{$T=296$} 
		{$T=428$} 
	\loigiai{ 
		Theo tính chất của cấp số cộng, ta có $x+z=2y$.\\ 
		Kết hợp với giả thiết $x+y+z=30$,  ta suy ra $3y=30\Leftrightarrow y=10$.\\
		Gọi $d$ là công sai  của cấp số cộng thì $x=y-d=10-d$ và $z=y+d=10+d$.\\
		$x+2$; $y+2$; $z+18$ là cấp số  nhân hay $12-d$, $12$, $28+d$.\\
		Theo tính chất  của cấp số nhân, ta có
		\begin{eqnarray*}
			 (12-d)(28+d)=12^2 &\Leftrightarrow & d^2+16d-192=0\\
			 &\Leftrightarrow & \hoac{&d=8\Rightarrow (x;y;z)=(2;10;18)\\&d=-24\Rightarrow (x;y;z)=(34;10;-14)&&(\text{loại}).}
		\end{eqnarray*}
		 Vậy $T=x^2+z^2=18^2+2^2=328$.
	 }
\end{ex}

% nguồn: Ok_Toan11_TracNghiem_Lan2
\begin{ex}%[1D2V3-7]%[ma:7d240a0b]
	Tính tổng  của cấp số nhân lùi vô hạn $(u_n)$ biết $u_1=1$ và $u_1$, $u_3$, $u_4$ theo thứ tự là ba số hạng liên tiếp trong một cấp số cộng.
	 \choice
		{$\dfrac{\sqrt{5}+1}{2}$} 
		{\True  $\dfrac{\sqrt{5}-1}{2}$}
		 {$\dfrac{1}{\sqrt{5}-1}$}
		{$2$} 
	\loigiai{ 
		$(u_n)$ là cấp số nhân lùi vô hạn có công bội $q$, suy ra $\left|q\right|< 1$ và $u_3=u_1\cdot q^2=q^2$, $u_4=u_1\cdot  q^3=q^3$.\\
		Mà và $u_1$, $u_3$, $u_4$ theo thứ tự là ba số hạng liên tiếp trong một cấp số cộng  nên $u_1+u_4=2u_3$.\\
		Từ đó ta có $1+q^3=2\cdot q^2\Leftrightarrow  q^3-2\cdot q^2+1=0\Leftrightarrow (q-1)(q^2-q-1)=0\Leftrightarrow q^2-q-1=0$
		$\Leftrightarrow  \hoac{&q=\dfrac{1+\sqrt{5}}{2}\\&q=\dfrac{1-\sqrt{5}}{2}}\Rightarrow q=\dfrac{1-\sqrt{5}}{2}$.\\
		 Vậy $S=\dfrac{u_1}{1-q}=\dfrac{1}{1-\dfrac{1-\sqrt{5}}{2}}=\dfrac{2}{1+\sqrt{5}}=\dfrac{\sqrt{5}-1}{2}$.
	} 
\end{ex}


%%%%%%%%%%%%%%%% ĐÚNG/SAI %%%%%%%%%%%%%%%%

% nguồn: Ok_dungsai_lop11_Lan1
\begin{ex}%[1D2K3-7]%[ma:28ab8458]
 Hãy cho biết tính đúng sai của mỗi mệnh đề  dưới đây?
 \choiceTF 
  {Cho hai số dương $ a $ và $ b $ không vượt quá $ 10 $ sao cho $ a-b $; $ 2 $; $ b $ theo thứ tự tạo thành một cấp số cộng và $ a+b $; $ 3a-2b $; $ 5a $  theo thứ tự lập thành một cấp số nhân. Khi đó $ a+b=8 $}
  {\True Cho ba  số $a$, $b$, $c$ là ba số liên tiếp của một cấp số cộng có công sai là $2$. Nếu tăng số thứ nhất thêm $1$, tăng số thứ hai thêm $1$ và tăng số thứ ba thêm $3$ thì được ba số mới là ba số liên tiếp của một cấp số nhân. Khi đó $a+b+c=9$}
  {\True Các số nguyên dương $x$, $y$ thỏa mãn ba số $x$; $2y$; $2x+3y-1$ theo thứ tự lập thành một cấp  số cộng và ba số $x$; $y-1$; $8$ theo thứ tự lập thành một cấp số nhân. Khi đó $x^2+2y$ bằng $14$}
  {Ba số $x$, $y$, $z$ theo thứ tự đó lập thành một cấp số nhân có công bội khác  $1$. Đồng thời chúng lần lượt là số hạng đầu, số hạng thứ hai và số hạng thứ năm của một cấp số cộng, biết rằng tống của chúng là $26$. Tích của ba số đó bằng $264$}
  \loigiai{
  \begin{itemchoice}
   \itemch Vì $ a-b $; $ 2 $; $ b $ theo thứ tự tạo thành một  cấp số cộng nên $ 2-(a-b)=b-2 \Leftrightarrow a=4 $.\\
   Vì $ a+b $; $ 3a-2b $; $  5a $ theo thứ tự lập thành một cấp số nhân nên $$ \dfrac{3a-2b}{a+b}=\dfrac{5a}{3a-2b}\Leftrightarrow \dfrac{12-2b}{4+b}=\dfrac{20}{12-2b}\Leftrightarrow \hoac{&b=1\\&b=16 \text{ (loại).}} $$
   Vậy  $ S=a+b=4+1=5. $
   \itemch Do $a,b,c$ là ba số liên tiếp của một cấp số  cộng có công sai là 2 nên $b=a+2, c=a+4$. \\
   $a+1, a+3,a+7$ là ba số liên tiếp của một cấp số nhân $ \Leftrightarrow (a+1)(a+7)=(a+3)^2 \Leftrightarrow  a=1$. \\
   Với $a=1$, ta có $\left\{\begin{aligned}
   &b=3\\ 
   &c=5\\ 
  \end{aligned}\right.. $ \\
  Suy ra $a+b+c=9$. 
  \itemch  Áp dụng tính chất cấp số cộng và cấp số nhân ta có hệ phương trình
   {\allowdisplaybreaks
  \begin{eqnarray*}
   &&\heva{&x+(2x+3y-1)=4y\\&8x=(y-1)^2} 
   \Leftrightarrow 
   \heva{&y=3x-1\\&8x=(3x-2)^2}\\ 
    &\Leftrightarrow&
   \heva{&y=3x-1\\&9x^2-20x+4=0} 
   \Leftrightarrow 
   \hoac{&\heva{&x=\dfrac{2}{9}\\&y=-\dfrac{1}{3}\text{  (loại)} }\\&\heva{&x=2\\&y=5.}}
  \end{eqnarray*}}
  Khi $x=2$ và $y=5$  thì $x^2+2y=14$.
  \itemch Do $x$, $y$, $z$ theo thứ tự đó lập thành một cấp số nhân có công bội  khác $1$ nên $x\ne y\ne z$.\\
  Do $x$, $y$, $z$  theo thứ tự đó là số hạng đầu, số hạng thứ hai và số hạng thứ năm của một cấp số cộng nên ta gọi $y=x+d$, $z=x+4d$ ($d\ne 0$).\\
  Theo  bài ra ta có
  $$\heva{&x(x+4d)=(x+d)^2\\  &3x+5d=26}\Leftrightarrow\heva{&2dx=d^2\\ &3x+5d=26}\Leftrightarrow\heva{&2x=d\\ &3x+5d=26}\Leftrightarrow \heva{&x=2\\ &d=4.}$$
  Suy  ra $x=2$, $y=6$ và $z=18$. Vậy $x\cdot y\cdot z=216$.
 \end{itemchoice}
  }
\end{ex}


%%%%%%%%%%%%%%%% TRẢ LỜI NGẮN %%%%%%%%%%%%%%%%

% nguồn: Ok_TraLoiNgan_Lop11_Lan2
\begin{ex}%[1D2H3-7]%[ma:b1f8689e]
	Số hạng thứ hai, số hạng đầu và số hạng thứ ba của một cấp số cộng với công sai khác $0$ theo thứ tự đó lập thành một cấp số  nhân với công bội $q$. Tìm $q$.
	\shortans{-2} 
	 \loigiai{
		Giả sử ba số hạng $a$, $b$, $c$ lập  thành cấp số cộng thỏa yêu cầu, khi đó $b$, $a$, $c$ theo thứ tự đó lập thành cấp số nhân công bội $q$. Ta có
		 \[
		\heva{&a  + c = 2b\\ &a = bq\\ &c = bq^2}
		\Rightarrow bq + bq^2 = 2b \Leftrightarrow \hoac{&b = 0\\ &q^2 + q - 2 = 0.} 
		 \]
		Nếu $b =  0 \Rightarrow a = b = c = 0$ nên $a$, $b$, $c$ là cấp số cộng công sai $d = 0$ (vô lí).\\
		Nếu $q^2 + q - 2 = 0 \Leftrightarrow q = 1$ hoặc $q = -2$. Nếu $q = 1 \Rightarrow a = b = c$  (vô lí), do đó $q = -2$.
	 }
\end{ex}

% nguồn: Ok_TraLoiNgan_Lop11_Lan2
\begin{ex}%[1D2H3-7]%[ma:8c5ec8cd]
	Các số $x + 6y$, $5x + 2y$, $8x + y$  theo thứ tự đó lập thành một cấp số cộng; đồng thời các số $x - 1$, $y + 2$, $x - 3y$ theo thứ tự đó lập thành một cấp số nhân. Tính $x^2 + y^2$.
	\shortans{40} 
	\loigiai{ 
		Theo giả thiết ta có $\heva{&(x + 6y) + (8x + y) = 2(5x +  2y)\\ &(x - 1)(x - 3y) = (y + 2)^2}$
		 \[
		 \Leftrightarrow
		\heva{&x  = 3y\\ &(3y - 1)(3y - 3y) = (y + 2)^2}
		 \Leftrightarrow
		\heva{&x = 3y\\ &0 = (y +  2)^2}
		\Leftrightarrow 
		\heva{&x = -6\\ &y  = -2.}
		 \]
		Suy ra $x^2 + y^2 =  40$.
	} 
\end{ex}

% nguồn: Ok_TraLoiNgan_Lop11_Lan1
\begin{ex}%[1D2V3-7]%[ma:1862198e]
	Quốc gia $X$ có trữ lượng dầu mỏ là $120$ triệu thùng vào đầu năm $2022$. Theo dự báo với mức khai thác dầu không đổi như năm $2022$ thì trữ lượng dầu của nước này sẽ hết sau đúng $50$ năm nữa. Do nhu cầu phát triển kinh tế, kể từ năm $2023$ trở đi mức khai thác dầu mỏ mỗi năm của quốc gia $X$ tăng lên $5 \%$ so với năm trước. Hỏi đến hết năm $2035$ thì  trữ lượng dầu của nước này còn lại bao nhiêu thùng (làm tròn đến hàng đơn vị)?
	 \shortans{73}
	 \loigiai{
		Trữ lượng dầu mỏ khai  thác trong năm $2022$ là $\dfrac{120}{50}=\dfrac{12}{5}$ triệu thùng.\\
		Trữ lượng khai thác dầu mỏ hàng năm  tính từ năm $2022$ đến năm $2035$ là
		 $$\dfrac{12}{5}; \dfrac{12}{5}\cdot 1{,}05;\dfrac{12}{5}\cdot \left(1{,}05\right)^2;\dfrac{12}{5}\cdot \left(1{,}05\right)^3;\cdots;\dfrac{12}{5}\cdot \left(1{,}05\right)^{13}.$$
		Dãy trên là một cấp số nhân có $u_1=\dfrac{12}{5}$,  công bội $q=1{,}05$.\\
		Tổng trữ lượng dầu  khai thác từ năm $2022$ đến năm $2035$ là
		$$S_{14}=u_1\cdot \dfrac{1-q^{14}}{1-q}=\dfrac{12}{5}\cdot \dfrac{1-\left(1{,}05\right)^{14}}{1-1{,}05}\approx 47 \text{ (triệu  thùng).}$$
		Vậy tính đến hết năm $2035$ thì trữ lượng dầu của nước này  còn lại $120-47=73$ triệu thùng.
	} 
\end{ex}

% nguồn: Ok_TraLoiNgan_Lop11_Lan2
\begin{ex}%[1D2V3-7]%[ma:ae28f6bc]
	Cho dãy số tăng $a$, $b$, $c$ ($c \in \mathbb{Z}$) theo thứ tự lập thành cấp số nhân; đồng thời $a$, $b + 8$, $c$ theo thứ tự lập thành cấp số cộng và $a$, $b + 8$, $c + 64$ theo thứ tự lập thành cấp số nhân. Tính giá trị biểu thức $P = a  - b + 2c$.
	\shortans{64} 
	 \loigiai{
		Ta  có
		\[ 
		\heva{&ac = b^2\\ &a + c = 2(b + 8)\\ &a(c + 64) = (b  + 8)^2}
		 \Leftrightarrow
		\heva{&ac = b^2 \quad \text{(1)}\\ &a - 2b = 16 - c \quad \text{(2)}\\  &ac + 64a = (b + 8)^2. \quad \text{(3)}}
		\] 
		Thay (1)  vào (3) ta được:
		\[ 
		b^2 + 64a = b^2 + 16b + 64 \Leftrightarrow 4a - b = 4.  \quad \text{(4)}
		\] 
		Kết hợp (2) với (4) ta được: 
		\[ 
		\heva{&a - 2b = 16 - c\\ &4a - b = 4} \Leftrightarrow \heva{&a = \dfrac{c - 8}{7}\\ &b  = \dfrac{4c - 60}{7}.} \quad \text{(5)}
		\] 
		Thay (5)  vào (1) ta được:
		\[ 
		7(c - 8)c = (4c - 60)^2 \Leftrightarrow  9c^2 - 424c + 3600 = 0
		\Leftrightarrow \hoac{&c = 36\\ &c = \dfrac{100}{9}} \Rightarrow c = 36  \ (\text{vì } c \in \mathbb{Z}).
		\] 
		Với  $c = 36 \Rightarrow a = 4$, $b = 12 \Rightarrow P = 4 - 12 + 72 = 64$.
	 }
\end{ex}

% nguồn: Ok_TraLoiNgan_Lop11_Lan1
\begin{ex}%[1D2C3-7]%[ma:44dbfee4]
	Cho ba số $a$, $b$, $c$ phân biệt có tổng bằng $49$. Ba số này theo thứ tự lập thành cấp số nhân và có thể coi chúng lần lượt là số hạng thứ $2$, thứ $5$ và thứ $10$ của một  cấp số cộng. Tính giá trị biểu thức $T=a^2+b^2+c^2$.
	 \shortans{931}
	\loigiai{ 
		Theo  đề bài, ta có\\
		Tổng ba số  $a$, $b$, $c$ bằng $49$, suy ra $a+b+c=49$. \tagEX{1}\noindent
		Ba số $a$, $b$, $c$ lập thành một cấp số  nhân, suy ra $b^2=a\cdot c$.\tagEX{2}\noindent
		Ba số $a$,  $b$, $c$ lần lượt là số hạng thứ $2$, thứ $5$ và thứ $10$ của một cấp số cộng, suy ra $\heva{&a=u_1+d\\&b=u_1+4d\\&c=u_1+9d.}$\tagEX{3}\noindent
		Từ $(1)$, $(2)$ và  $(3)$ suy ra $\heva{&3u_1+14d=49\\&\left(u_1+4d\right)^2=\left(u_1+d\right)\left(u_1+9d\right).~(*)}$\\
		Phương trình $(*)  \Leftrightarrow 7d^2-2u_1d=0 \Leftrightarrow \hoac{&d=0\\&u_1=\dfrac{7}{2}d.}$\\
		Với $d=0$, suy ra $a=b=c$ (loại  do $a$, $b$, $c$ phân biệt).\\
		Với  $u_1=\dfrac{7}{2}d $, ta có
		\[3u_1+14d=49 \Rightarrow 3\cdot \dfrac{7}{2}d+14d=49 \Rightarrow d=2  \Rightarrow u_1=7.\]
		Suy  ra $\heva{&a=9\\&b=15\\&c=25}\Rightarrow T=a^2+b^2+c^2=931$.
	 }
\end{ex}

% nguồn: Ok_TraLoiNgan_Lop11_Lan1
\begin{ex}%[1D2C3-7]%[ma:f4f0c177]
	Một quả bóng được thả thẳng đứng từ độ cao $8$ mét rơi xuống đất và nảy lên. Giả sử sau mỗi một lần rơi xuống, nó nảy lên được một độ cao  bằng $60\%$ độ cao vừa rơi xuống.
	Tính tổng quãng đường quả bóng di chuyển được kể từ lúc thả xuống đến khi quả bóng  chạm đất lần thứ 8 (làm tròn kết quả đến hàng phần mười của mét).\shortans{19,7}
	\loigiai{Gọi $u_n$ (m) là độ cao mà quả bóng đạt được sau  khi nảy lên ở lần thứ $n$.\\
		Khi đó, dãy $\left(u_n\right)$ lập thành cấp số nhân có $u_1 = 8 \cdot 0{,}6 = 4{,}8$ và công  bội $q = 0{,}6$. \\
		Kể từ lúc thả xuống đến khi quả bóng chạm đất lần thứ $8$, quả bóng đã được nảy lên $7$ lần rồi lại rơi xuống. Do quãng đường quả bóng nảy lên và rơi xuống bằng nhau nên tổng quãng đường quả bóng di chuyển được kể  từ lúc thả xuống đến khi quả bóng chạm đất lần thứ $8$ là
		$$ S= 8 + 2 \left(u_1 + u_2 + \ldots +  u_{7}\right) = 8 + 2 \cdot 4{,}8 \cdot \dfrac{1- 0{,}6^{7}}{1- 0{,}6} \approx 19{,}7 \,(\mathrm{m}).$$
	 }
\end{ex}


%%%%%%%%%%%%%%%% TỰ LUẬN %%%%%%%%%%%%%%%%

% nguồn: Ok_TuanLuan_Toan11_Lan2
\begin{ex}%[1D2H3-7]%[ma:1854ad15]
	Tìm $a$, $b$ biết rằng $1$, $a$, $b$ là $3$ số hạng liên tiếp của cấp số cộng và $1$, $a^2$, $b^2$ là $3$  số hạng liên tiếp của một cấp số nhân.
	 \loigiai{
		Theo đề bài  ta có hệ phương trình $\heva{&1+b=2a\\&b^2=a^4} \Leftrightarrow \heva{&1+b=2a~\quad(1)\\&b=\pm a^2.}$\\
		Với $b=a^2$ thay vào $(1)$ được $1+a^2 = 2a \Leftrightarrow a^2-2a+1=0  \Leftrightarrow a=1 \Rightarrow b=1$.\\
		Với $b=-a^2$ thay vào $(1)$ được $1-a^2 = 2a \Leftrightarrow -a^2-2a+1=0 \Leftrightarrow  \hoac{&a=-1+\sqrt{2}\\&a=-1-\sqrt{2}.}$
		\begin{itemize}
			\item $a=-1+\sqrt{2} \Rightarrow b = -\left(-1+\sqrt{2}\right)^2 \Leftrightarrow b=-3+2\sqrt{2}$. 
			\item $a=-1-\sqrt{2} \Rightarrow b  = -\left(-1-\sqrt{2}\right)^2 \Leftrightarrow b=-3-2\sqrt{2}$.
		\end{itemize}
		Kết  luận $\heva{&a=1\\&b=1}$; $\heva{&a=-1+\sqrt{2}\\&b=-3+2\sqrt{2}}$; $\heva{&a=-1-\sqrt{2}\\&b=-3-2\sqrt{2}}$ thỏa mãn yêu cầu đề bài.
	 }
\end{ex}
